\section{Kluczowe komponenty kodu}
\label{sec:kod}

Cała logika modelu jest wydzielona do \texttt{src/model/mlp.py} i pokryta testami
jednostkowymi (\texttt{tests/model/}). Poniżej omówiono najistotniejsze fragmenty.

\subsection{\texttt{TfidfDataset}: leniwa densyfikacja}

Macierz TF-IDF jest rzadka (50\,000 kolumn, lecz nieliczne niezerowe wartości w
wierszu). Gęste zmaterializowanie całego korpusu naraz byłoby kosztowne pamięciowo,
dlatego \texttt{TfidfDataset} przechowuje dane jako macierz rzadką i densyfikuje
\emph{pojedynczy wiersz dopiero przy pobraniu} przez \texttt{DataLoader}.

\begin{lstlisting}[caption={Leniwa densyfikacja wiersza w \texttt{TfidfDataset}.}]
def __getitem__(self, index: int) -> tuple[torch.Tensor, torch.Tensor]:
    row = self.features[index].toarray().ravel().astype(np.float32, copy=False)
    return torch.from_numpy(row), self.labels[index]
\end{lstlisting}

\subsection{\texttt{MLPClassifier}: sieć zwracająca logity}

Sieć budowana jest jako stos bloków \texttt{Linear $\to$ ReLU $\to$ Dropout},
zakończony pojedynczą warstwą liniową. \texttt{forward} zwraca surowe \emph{logity}
(bez sigmoidy), co pozwala użyć numerycznie stabilnej \texttt{BCEWithLogitsLoss}
w treningu, a \texttt{torch.sigmoid} dopiero przy wnioskowaniu.

\begin{lstlisting}[caption={Konstrukcja warstw \texttt{MLPClassifier}.}]
layers: list[nn.Module] = []
in_features = input_dim
for width in hidden_dims:
    layers.append(nn.Linear(in_features, width))
    layers.append(nn.ReLU())
    layers.append(nn.Dropout(dropout))
    in_features = width
layers.append(nn.Linear(in_features, 1))
self.net = nn.Sequential(*layers)
\end{lstlisting}

\subsection{\texttt{train\_mlp}: pętla treningowa z early stopping}

Pętla naprzemiennie trenuje (\texttt{\_train\_one\_epoch}) i waliduje
(\texttt{\_validate}, zwraca stratę i F1). Po każdej epoce, jeśli F1 walidacji
poprawił się, zapisywany jest stan wag; w przeciwnym razie zliczane są epoki bez
poprawy, a po przekroczeniu \texttt{patience} trening zostaje przerwany. Na końcu
przywracane są najlepsze wagi.

\begin{lstlisting}[caption={Rdzeń early stopping w \texttt{train\_mlp}.}]
optimizer = torch.optim.Adam(model.parameters(), lr=lr, weight_decay=weight_decay)
criterion = nn.BCEWithLogitsLoss()
...
for _ in range(epochs):
    train_loss = _train_one_epoch(model, train_loader, optimizer, criterion, dev)
    val_loss, val_f1 = _validate(model, val_loader, criterion, dev)
    ...
    if val_f1 > best_val_f1:
        best_val_f1 = val_f1
        best_state = {n: t.detach().cpu().clone()
                      for n, t in model.state_dict().items()}
        epochs_without_improvement = 0
    else:
        epochs_without_improvement += 1
        if epochs_without_improvement >= patience:
            break
if best_state:
    model.load_state_dict(best_state)
\end{lstlisting}

\subsection{\texttt{predict\_proba} i \texttt{evaluate\_mlp}}

\texttt{predict\_proba} przepuszcza dane przez sieć w trybie ewaluacji i zwraca
$P(\text{Real})$ jako tablicę \texttt{numpy} (logit $\to$ \texttt{sigmoid}).
\texttt{evaluate\_mlp} owija to w obliczenie metryk i zwraca \emph{dokładnie taki
sam słownik} (\texttt{accuracy}, \texttt{precision}, \texttt{recall}, \texttt{f1},
\texttt{roc\_auc}, \texttt{confusion\_matrix}) co \texttt{baseline.evaluate}, dzięki
czemu porównanie obu modeli jest bezpośrednie.

\subsection{Reprodukowalność i urządzenie}

\texttt{set\_seed} ustawia ziarno dla NumPy oraz PyTorcha (CPU i CUDA/ROCm), a
\texttt{resolve\_device} automatycznie wybiera GPU, gdy jest dostępne (na maszynie
deweloperskiej działa build ROCm, dla którego \texttt{torch.cuda.is\_available()}
zwraca \texttt{True}). Oba modele korzystają z tego samego wektoryzatora
(\texttt{baseline.make\_tfidf\_vectorizer}), co gwarantuje identyczną reprezentację
tekstu i uczciwe porównanie architektur.
